% Original demonstration assembly; official template files remain unchanged.
\documentclass[degree=master,fontset=fandol]{thuthesis}
\input{thusetup}
\thusetup{
  output=electronic,
  title={学位论文写作与自查指南},
  title*={A Guide to Thesis Writing and Self-Review},
  author={ThuThesis Authoring 示例项目},
  author*={ThuThesis Authoring Example Project},
  department={}, discipline={}, discipline*={},
  degree-category={}, degree-category*={},
  supervisor={}, supervisor*={}, associate-supervisor={}, associate-supervisor*={},
}
\hypersetup{pdftitle={学位论文写作与自查指南},pdfauthor={ThuThesis Authoring Example Project}}
\begin{document}
\begin{titlepage}
\centering
\vspace*{35mm}
{\zihao{1}\bfseries 学位论文写作与自查指南\par}
\vspace{12mm}
{\large A Guide to Thesis Writing and Self-Review\par}
\vspace{20mm}
{\zihao{3} 原创教学示例\par}
\vspace{8mm}
\begin{minipage}{.78\textwidth}
\centering
本文用于展示材料组织、LaTeX 排版和质量检查。\\
不是学校官方规范，不是可提交的正式学位论文。\\
数值 2、4、6 仅为教学数据，不代表真实实验。\\
两项未确认要求以 TODO 保留。
\end{minipage}
\vfill
ThuThesis Authoring 示例项目\par
2026\par
\end{titlepage}
\frontmatter
% !TEX root = ../main.tex
\begin{abstract}
学位论文写作需要把研究问题、材料、方法和结论组织成可追溯的论证。本指南围绕这一过程，介绍章节规划、段落组织、公式与图表表达、引用管理以及提交前自查。指南通过原创对照句说明如何将含糊表述改为边界明确的文字，以三个教学数值展示变量定义和均值计算，并给出逐步修订的流程。对于院系要求和导师意见等尚未确认的信息，指南保留明确的待办标记。示例旨在帮助读者理解可编辑论文工程的组织方式，以及内容审阅、结构检查和真实编译各自能够提供的证据。它不以编译成功代替学术判断，也不把通用建议视为具体学校的提交规范。
\thusetup{keywords={学位论文, 写作方法, 证据管理, 质量检查}}
\end{abstract}

\begin{abstract*}
Thesis writing requires a traceable argument connecting research questions, source material, methods, and conclusions. This guide introduces chapter planning, paragraph development, equations and tables, citation management, and checks before submission. Original sentence pairs illustrate how vague statements can be replaced with bounded descriptions. Three illustrative values demonstrate variable definitions and mean calculation, while a workflow explains iterative revision. Unconfirmed departmental requirements and supervisory preferences remain explicit action items. The example explains the organization of an editable thesis project and the distinct evidence provided by content review, structural checks, and actual compilation. A successful build does not replace academic judgment, and general writing advice does not establish institution-specific submission requirements.
\thusetup{keywords*={thesis writing, writing methods, evidence management, quality checks}}
\end{abstract*}

\tableofcontents
\mainmatter
% !TEX root = ../main.tex
\chapter{从研究问题到论文结构}
\label{chap:guide1}

\section{先写清楚问题，再安排章节}

开始写作时，可以先用一小段话说明研究对象、希望回答的问题、已有材料和当前限制。问题应当具体到能够判断哪些材料与它有关。例如，“研究写作工具”只是主题；“说明一份材料如何组织为可编辑章节，并记录仍需确认的信息”才给出了可执行的写作目标。这里的区别不是句子长短，而是读者能否据此理解讨论范围。尚未取得的材料应单独列出，不要把预计完成的工作写成已经得到的结果。

问题陈述还需要与研究结论相互对应。如果开头讨论的是引用信息如何追溯，结尾就应回答来源是否明确、哪些信息仍不完整，而不能突然转向工具是否提升科研创新。写作过程中可以修改问题，但每次调整都要检查摘要、章节标题和结论是否同步。为同一个概念保持相同名称，通常比增加更多术语更有帮助。本项目配套的章节组织说明给出了问题、材料和章节之间的记录方法。\cite{guide1}

\section{为材料建立可追溯的清单}

在正文之外维护一份材料清单，记录每项材料的名称、来源、版本、准备使用的位置和待核对事项。自己的记录也应注明产生时间与用途；引用他人材料时，应保留能够定位来源的信息。材料清单的目的不是积累文件数量，而是让每个论点都能回到具体证据。如果两个文件来自不同修改阶段，先判断哪一个是本次写作的依据，再合并内容，避免把相互矛盾的数值拼在同一段中。

可以将“材料已经提供”“材料已经核对”和“内容已经写入”作为不同状态。收到一张表格，并不等于表中指标的定义已经清楚；把文字放进章节，也不等于完成了对照审阅。当材料缺少单位、条件或来源时，应记录缺项与影响范围。这里不要求使用专门的软件，一个包含固定字段的普通文本表格就可以承担这些工作。涉及未公开数据时，公开示例应另用独立的教学材料。

\section{用章节职责约束写作范围}

章节提纲需要说明每一章解决什么问题，而不只是列出几个宽泛标题。表~\ref{tab:guide1} 给出一种通用的职责划分；实际论文应根据研究类型调整，不必机械套用。理论研究、实验研究和工程实践的材料不同，章节之间的关系也可能不同。判断提纲是否清楚，可以尝试用一句话描述相邻两章的衔接：前一章给出了什么，后一章为什么需要接着讨论。

\begin{table}[htbp]
  \centering
  \small
  \renewcommand{\arraystretch}{1.25}
  \caption{常见章节职责与检查问题}
  \label{tab:guide1}
  \begin{tabular}{@{}>{\raggedright\arraybackslash}p{.2\linewidth}>{\raggedright\arraybackslash}p{.52\linewidth}>{\raggedright\arraybackslash}p{.2\linewidth}@{}}
    \toprule
    章节角色 & 主要职责 & 自查问题 \\
    \midrule
    绪论 & 明确问题、背景与范围 & 读者能否说出本文要回答什么？ \\
    方法 & 交代步骤、变量与适用条件 & 信息是否足以理解操作过程？ \\
    结果与讨论 & 展示观察并解释其边界 & 结论能否回到具体材料？ \\
    总结 & 回答问题并列出局限 & 是否引入了正文未说明的新结论？ \\
    \bottomrule
  \end{tabular}
\end{table}

摘要应在正文结构稳定后回头校对，至少交代讨论对象、采用的方式、主要内容和适用边界。中英文摘要可以采用各自自然的表达，但研究对象、数量、关键条件和结论强度应一致。章节规划并不意味着从第一章顺序写到最后一章；作者可以先整理证据最充分的部分，再补上连接这些部分的解释。修改顺序可以灵活，论证关系必须清楚。

% !TEX root = ../main.tex
\chapter{正文、图表与引用的写法}
\label{chap:guide2}

\section{一段文字只推进一个主要判断}

段落可以按照“主要判断、相关材料、解释、边界”的顺序组织。主要判断告诉读者这一段要说明什么，材料给出依据，解释连接材料与判断，边界则说明哪些情况尚不能推出。并非每一段都必须包含四个句子，但这四个问题可以用来发现论证中的缺口。把多个独立判断塞进同一长句，往往使读者难以确定哪些证据支持哪个结论。

修改文字时，应优先修正含义和证据关系，再处理语气与措辞。将“效果很好”替换为“效果显著”，只是换了一个词，并没有补足指标、条件和依据。表~\ref{tab:guide2} 中的推荐写法保留了未知信息，没有把待完成的检查改写为确定结果。它们是原创教学对照句，不是从实际论文摘录的文字。证据与引用清单可以辅助记录每次修改是否改变了原有结论的强度。\cite{guide2}

\begin{table}[htbp]
  \centering
  \small
  \renewcommand{\arraystretch}{1.25}
  \caption{表述修改的教学对照}
  \label{tab:guide2}
  \begin{tabular}{@{}>{\raggedright\arraybackslash}p{.2\linewidth}>{\raggedright\arraybackslash}p{.52\linewidth}>{\raggedright\arraybackslash}p{.2\linewidth}@{}}
    \toprule
    原始表述 & 推荐表述 & 修改目的 \\
    \midrule
    这个方法效果很好。 & 需要先说明评价指标与比较条件，再描述观察到的变化。 & 补足判断依据 \\
    所有情况都适用。 & 本文仅讨论材料中明确列出的条件，其他情形尚未验证。 & 限定适用范围 \\
    参考文献已经整理完了。 & 已完成条目录入，作者、年份与正文引用仍需逐项核对。 & 区分操作与验证 \\
    \bottomrule
  \end{tabular}
\end{table}

\section{公式需要定义、计算和文字解释}

公式进入正文之前，先说明它回答什么问题，并在首次使用时定义符号。以式~\eqref{eq:mean} 的算术平均数为例，设 $x_i$ 表示第 $i$ 个教学数值，$n$ 表示数值个数，$\bar{x}$ 表示算术平均数，则

\begin{equation}
\bar{x}=\frac{1}{n}\sum_{i=1}^{n}x_i.
\label{eq:mean}
\end{equation}

本例的数据为 2、4、6，均为独立构造的教学数值，单位为任意教学单位。于是 $n=3$，数值之和为 12，均值为 4。这个计算仅说明公式、变量和表述如何对应，不代表任何实际实验的测量结果。正文引用公式时应使用自动编号，不要手工写死编号。公式后的文字还应解释结果能够回答的问题，避免读者只看到一个孤立数值。

同一个符号应尽量保持同一含义。如果后文需要改变单位、集合或统计范围，必须在改变的位置说明。不要仅凭均值判断数据的离散程度，也不要从三个教学数值推导方法的一般有效性。这里保留简单算例，是为了让读者能够独立核算，并检查正文、公式和数据文件是否一致。复杂研究仍需根据实际设计选择指标和报告方式。

\section{图表与引用都应能够回到来源}

表格适合并列比较，图形适合呈现关系或流程。选择形式时，先判断读者需要完成什么阅读任务，再决定是否增加图表。标题应说明图表的对象，表头应给出字段含义，涉及数值时还要交代单位和统计范围。正文需要指出读者应观察什么，而不是只写“如下图所示”。修改图表顺序后，应重新检查标题、标签和正文引用是否仍然对应。

引用用于说明某个定义、方法或判断来自哪里，不能替代作者自己的解释。将条目放进文献数据库之后，仍需检查正文是否确实引用了它，引用位置是否符合所支撑的内容。对于缺失的作者、年份或出版信息，应回到材料核对；无法确认的字段保持缺失，不凭印象补齐。本指南引用的两份说明均随项目提供，属于原创教学文档，并不声称是外部发表的研究成果。

% !TEX root = ../main.tex
\chapter{修改与提交前自查}
\label{chap:guide3}

\section{从内容关系开始安排修改顺序}

修改可以从影响范围较大的问题开始。先检查研究问题是否明确，章节是否共同回答这个问题，再检查段落之间的连接，最后检查措辞、标点与版式。若论证结构仍在变化，过早逐字调整格式可能带来重复工作。每轮修改可以列出一个明确目标，例如检查所有结论是否有对应材料，或者统一变量名称，而不是笼统地要求把全文变得更好。

图~\ref{fig:workflow} 表现“材料整理、初稿、检查、修改、编译”这五个步骤。检查发现问题后，可以回到相关材料或章节重新处理；流程中的编译负责验证当前工程能否产生排版结果，不承担学术评价。把修改原因和涉及文件记下来，有助于后续解释某段内容为什么发生变化，也便于区分本轮新出现的问题与此前已经记录的待办。

\begin{figure}[htbp]
\centering
\fbox{\parbox[c][12mm][c]{.13\linewidth}{\centering 材料\\整理}}\hfill
$\rightarrow$\hfill
\fbox{\parbox[c][12mm][c]{.13\linewidth}{\centering 初稿}}\hfill
$\rightarrow$\hfill
\fbox{\parbox[c][12mm][c]{.13\linewidth}{\centering 检查}}\hfill
$\rightarrow$\hfill
\fbox{\parbox[c][12mm][c]{.13\linewidth}{\centering 修改}}\hfill
$\rightarrow$\hfill
\fbox{\parbox[c][12mm][c]{.13\linewidth}{\centering 编译}}
\caption{从材料整理到编译的示例写作流程}
\label{fig:workflow}
\end{figure}


\section{将待确认事项保留在可检索的位置}

待办项应说明缺少什么、需要向谁核对以及影响哪些内容。只写“待完善”难以指导下一步行动；写明“核对提交要求中的摘要长度与附件清单”则更容易执行。待办既可以出现在相关段落旁，也可以汇总到交付说明中，但两处内容应保持一致。不能因为暂时无法获取答案，就用一个看似合理的数字或规则代替它。

本指南面向通用写作过程，尚未取得具体院系的最新要求和具体导师对章节安排的意见。因此保留以下两项待办：

\begin{itemize}
\item {}[TODO: 核对所在院系的最新提交要求。]
\item {}[TODO: 确认导师对章节安排的具体要求。]
\end{itemize}

完成这些待办需要读者结合自己的实际情况提供信息。公开教学示例保留它们，是为了说明缺项如何被发现和追踪，并不意味着所有论文都应使用同一份检查清单。确认后，应把依据、修改位置和结果一起记录，再运行与修改内容相关的检查。待办消失本身不能证明答案可靠，来源仍然需要审阅。

\section{分别报告内容审阅、结构检查和编译}

提交前可以从三个方面检查当前成果。内容审阅关注论点、证据和表述是否一致；结构检查关注文件路径、章节连接、标签与引用等是否完整；真实编译关注工具是否成功生成本次排版产物。三者相互补充，但不能相互替代。检查工具报告结构通过时，作者仍应阅读正文；生成 PDF 后，也应逐页观察公式、图表和段落是否发生裁切或错位。

对检查结果的记录应包含本次使用的材料版本、检查范围、实际执行的命令、发现的问题和未验证的部分。如果仍有待办，应明确说明哪些内容尚待确认。重复编译时，要确认看到的是本次生成的文件，而不是早先留下的 PDF。这样，交付者与接收者才能围绕同一组材料讨论，避免把文件存在、检查通过和正式提交混为一谈。

本指南最终交付的是可以继续编辑的示例工程。读者可用它理解章节组织与检查方法，再把这些方法应用到自己的材料。公开展示应突出原创正文、公式、表格和自查流程；真实姓名、学号以及未经确认的流程页面不应成为示例内容。具体论文的学术判断、院系规范和导师意见仍需由作者在实际写作中逐项落实。

\bibliography{ref/refs}
\end{document}
